
The optimality conditions for quadratic functions are related to the
positive definiteness of the (constant) second derivative matrix.
The derivatives of the objective function are
\[
\nabla f(x) = {Qx + g}, \; \nabla^2 f(x) = Q.
\]
The positive definiteness of $Q$ can be identified from its
eigenvalues, that are the roots of the characteristic
polynomial. Therefore, we solve the equation
\begin{align*}
\det(\lambda I - Q) = \det \begin{pmatrix}
1- \lambda & \alpha\\ \alpha & 1-\lambda
\end{pmatrix} & = (1-\lambda)^2 - \alpha^2\\
& = \lambda^2 - 2\lambda + (1-\alpha^2)\\
& = (1-\alpha - \lambda)(1+\alpha - \lambda).
\end{align*}
We obtain the eigenvalues:
\begin{align*}
\lambda_1 & = 1-\alpha\\
\lambda_2 & = 1+\alpha.\\
\end{align*}
The corresponding (normalized) eigenvectors are
\begin{align*}
u_1 &= \cvect{\frac{\sqrt{2}}{2} \\ -\frac{\sqrt{2}}{2}} & \text{ as } Qu_1 = \lambda_1 u_1,\\
u_2 &= \cvect{\frac{\sqrt{2}}{2} \\ \frac{\sqrt{2}}{2}} & \text{ as } Qu_2 = \lambda_2 u_2.
\end{align*}

We first investigate the optimality conditions.
\begin{description}
\item[Necessary conditions]  When $Q$ is not positive
  semidefinite, that is, when at least one eigenvalue is non
  positive, there is no optimal solution. This is the case if $\alpha > 1$ (so that $\lambda_1 <
  0$), or $\alpha  < -1$ (so that $\lambda_2 < 0$).
\item[Sufficient conditions] When $Q$ is positive definite, the
  problem has a unique global minimum. This is the case if $-1 <
  \alpha < 1$, so that $\lambda_1 > 0$ and $\lambda_2 > 0$.
\end{description}
Now, what happens when $Q$ is positive  semidefinite (so that the
necessary conditions are verified) but not positive definite (so that
the sufficient conditions are not verified)?

 This is the case if $\alpha = 1$, so that
  $\lambda_1=0$ and $\lambda_2=2$, or if $\alpha = -1$, so that
  $\lambda_1=2$ and $\lambda_2=0$.

 Geometrically, we can decompose the space into two subspaces. In the
 subspace spanned by the eigenvectors of the positive eigenvalues, the
 function is strictly convex. In this subspace, there is a unique minimum. In the subspace spanned by the
 eigenvectors of the zero eigenvalues, the function is linear, as
 there is no curvature. Therefore, it is bounded only if it is
 constant. And, in that case, there is an infinite number of minima. 

 Algebraically, we use the Schur decomposition of $Q$:
\begin{align*}
Q & = U \Lambda U^T\\
& = \begin{pmatrix}
\frac{\sqrt{2}}{2} & \frac{\sqrt{2}}{2}\\ -\frac{\sqrt{2}}{2} & \frac{\sqrt{2}}{2}
\end{pmatrix} \begin{pmatrix}
1-\alpha & 0\\ 0& 1+\alpha
\end{pmatrix} \begin{pmatrix}
\frac{\sqrt{2}}{2} & -\frac{\sqrt{2}}{2}\\ \frac{\sqrt{2}}{2} & \frac{\sqrt{2}}{2}
\end{pmatrix},
\end{align*}
where $U$ is an orthogonal matrix composed of the eigenvectors of $Q$
organized in columns.

Let's use the matrix $U$ to decompose $x$ and $g$:
\[
U^Tx = x' = \cvect{x'_1 \\ x'_2 } \text{ and } U^Tg = g' = \cvect{g'_1
\\ g'_2},
\]
that is
\[
x' = \frac{\sqrt{2}}{2}\cvect{x_1-x_2 \\ x_1+x_2}\text{ and } g' = \frac{\sqrt{2}}{2}\cvect{g_1-g_2 \\ g_1+g_2}.
\]


We use the Schur decomposition, and the fact that the U matrix is
orthogonal, so that $UU^T=I$, to rewrite the objective function as follows:
\begin{align*}
  f(x) &= \frac{1}{2} x^T Q x + g^T x + c \\
  &= \frac{1}{2} x^T U \Lambda U^Tx + g^TUU^Tx + c.
\end{align*}
Using the decomposition of the vectors, we obtain the function in the
new variables
\begin{align*}
  \tilde{f}(x') &= \frac{1}{2} {x'}^T \Lambda x' + {g'}^Tx' + c \\
        &=\frac{1}{2} \lambda_1(x'_1)^2 + \frac{1}{2}
  \lambda_2(x'_2)^2 + g'_1 x'_1 + g'_2 x'_2.
\end{align*}
The gradient of this function is
\[
\nabla \tilde{f}(x') = \cvect{\lambda_1 x_1' + g'_1 \\ \lambda_2 x'_2 + g'_2}.
\]
We see that, for the gradient to be zero (which is a necessary
condition for optimality), we need
\begin{align*}
  x'_i &= \frac{-g'_i}{\lambda_i} & \text{ if } \lambda_i \neq 0, \\
  g'_i &= 0 & \text{ if } \lambda_i = 0. \\
\end{align*}

Now, we consider the two values of $\alpha$
identified above.

\begin{itemize}
\item If $\alpha=1$, then $\lambda_1=0$ and $\lambda_2 = 2$, the
  gradient is zero when
  \[
g'_1 = \frac{\sqrt{2}}{2}(g_1-g_2) = 0 \text{ and } x'_2 = \frac{\sqrt{2}}{2}(x_1+x_2)=\frac{-g'_2}{2}.
  \]
It means that the function must be such that $g_1=g_2$. If so, any
solution such that
\[
x_1 = -g_1 - x_2
\]
is a global minimum. It can easily be verified that the gradient of the function
\[
f(x) = Qx+g = \cvect{x_1+\alpha x_2 + g_1 \\ \alpha x_1 + x_2 + g_2}
\]
is indeed zero when $\alpha=1$, $x_1 = -g_1 - x_2$ and $g_1=g_2$, for
any value of $x_2$.

If $g_1 \neq g_2$, the problem is unbounded and
has no optimal solution.

\item If $\alpha=-1$, then $\lambda_1=2$ and $\lambda_2 = 0$, and we need
  \[
 x'_1 = \frac{\sqrt{2}}{2}(x_1-x_2)=\frac{-g'_1}{2}\text{ and } g'_2 = \frac{\sqrt{2}}{2}(g_1+g_2) = 0.
  \]
It means that the function must be such that $g_1=-g_2$. If so, any
solution such that
\[
x_1 = -g_1+x_2
\]
is a global minimum.
It can easily be verified that the gradient of the function
\[
f(x) = Qx+g = \cvect{x_1+\alpha x_2 + g_1 \\ \alpha x_1 + x_2 + g_2}
\]
is indeed zero when $\alpha=-1$, $x_1 = -g_1 + x_2$ and $g_1=-g_2$, for
any value of $x_2$.

If $g_1\neq-g_2$, the problem is unbounded and has no
optimal solution. 
\end{itemize}

We can now answer the questions. 

\begin{enumerate}
\item \emph{For which values of $\alpha$ $f$ does not have a local
minimum?}
  \begin{itemize}
    \item If $\alpha > 1$, so that $\lambda_1 <
      0$ (from the necessary conditions),
    \item if $\alpha  < -1$, so that $\lambda_2 < 0$ (from the
      necessary conditions),
    \item if $\alpha = 1$ when $g_1\neq g_2$ (from the discussion above), 
    \item if $\alpha = -1$ when $g_1\neq -g_2$ (from the discussion above).
  \end{itemize}

\item \emph{For which values of $\alpha$ $f$ has a unique global
  minimum?}  If $-1 < \alpha < 1$, so that $\lambda_1
  > 0$ and $\lambda_2 > 0$ (from the sufficient conditions).

\item \emph{What are the conditions on $g$ and $\alpha$ so that the
problem has an infinite number of global minima?}
  From the discussion above,
  \begin{itemize}
  \item if $\alpha=1$ when $g_1=g_2$, 
  \item if $\alpha=-1$ when $g_1=-g_2$.
  \end{itemize}
\end{enumerate}
